% Source: tex/chapters/ch06/05_ソフトウェアエンジニアリング運用計画.tex
\subsubsection{ソフトウェア容量管理}

容量管理とは、現在と将来の需要に対して、十分な処理能力、記憶容量、ネットワーク、データベース性能などを確保する活動である。利用者数、トランザクション数、データ量、業務量の変化を予測し、それを具体的な要求に変換して文書化する。

容量管理は、性能と容量の問題の中心となる。新規サービスや変更サービスを設計するとき、想定負荷を見積もり、必要な資源をモデル化する。容量計画には、現在の性能、将来要求、費用を伴う選択肢、SLA達成のための推奨策を含める。

\begin{pointbox}{例}
動画配信サービスで、通常時は一日10万人が利用するが、イベント配信時には同時接続が10倍になると予測される場合、容量管理ではピーク時のサーバ数、ネットワーク帯域、キャッシュ戦略、負荷分散、監視しきい値を計画する必要がある。
\end{pointbox}
